% preamble-exam-zh.tex — 共享 preamble
% 用于所有 exam-zh 模板
%
% 样式策略：
%   屏幕 → 语义彩色（蓝=关键想法，红=易错，绿=路线，灰=提示/教师备注）
%   打印 → 自动降级为黑白（通过 print mode 或 monochrome 选项）

\documentclass{exam-zh}

% ---------- 基础配置 ----------
\examsetup{
  page/size = a4paper,
  page/show-head = false,
  page/show-foot = false,
  sealline/show = false,
  font = times,
  math-font = xits,
}

\usepackage{tcolorbox}
\usepackage{needspace}
\tcbuselibrary{skins,breakable}

% ---------- 打印/屏幕颜色切换 ----------
% 用法：\documentclass[monochrome]{exam-zh} 即可切换为纯黑白
% 注意：不要使用 \if@xxx 放在 makeatother 外部；这里改成普通条件名。
\newif\ifeduprintcolor
\eduprintcolortrue

\makeatletter
\@ifclasswith{exam-zh}{monochrome}{\eduprintcolorfalse}{}
\makeatother

% ---------- 语义颜色定义 ----------
% 屏幕：有色彩区分；打印：降级为灰度
\ifeduprintcolor
  % --- 屏幕彩色模式 ---
  \definecolor{edu-blue}{RGB}{47, 82, 122}       % 关键想法
  \definecolor{edu-blue-bg}{RGB}{243, 246, 252}
  \definecolor{edu-red}{RGB}{180, 40, 40}         % 易错提醒
  \definecolor{edu-red-bg}{RGB}{255, 245, 240}
  \definecolor{edu-green}{RGB}{34, 100, 60}       % 路线图
  \definecolor{edu-green-bg}{RGB}{242, 250, 245}
  \definecolor{edu-orange}{RGB}{160, 90, 0}       % 训练目标
  \definecolor{edu-orange-bg}{RGB}{255, 248, 235}
  \definecolor{edu-gray}{RGB}{100, 100, 100}      % 提示/教师备注
  \definecolor{edu-gray-bg}{RGB}{248, 248, 248}
  \definecolor{edu-purple}{RGB}{90, 50, 120}      % 思考题
  \definecolor{edu-purple-bg}{RGB}{248, 244, 252}
\else
  % --- 打印黑白模式 ---
  \definecolor{edu-blue}{gray}{0.15}
  \definecolor{edu-blue-bg}{gray}{0.96}
  \definecolor{edu-red}{gray}{0.25}
  \definecolor{edu-red-bg}{gray}{0.97}
  \definecolor{edu-green}{gray}{0.20}
  \definecolor{edu-green-bg}{gray}{0.97}
  \definecolor{edu-orange}{gray}{0.25}
  \definecolor{edu-orange-bg}{gray}{0.98}
  \definecolor{edu-gray}{gray}{0.40}
  \definecolor{edu-gray-bg}{gray}{0.97}
  \definecolor{edu-purple}{gray}{0.30}
  \definecolor{edu-purple-bg}{gray}{0.98}
\fi

% ---------- 自定义教学环境 ----------
% 共性：enhanced + breakable（跨页不断裂）

% 原题卡片 — 强边框，突出显示
\newtcolorbox{problemcard}{
  enhanced, breakable,
  colback=edu-blue-bg,
  colframe=edu-blue,
  boxrule=1.2pt,
  arc=0pt,
  left=3mm, right=3mm, top=3mm, bottom=3mm,
}

% 解题路线图 — 左边线 + 浅底
\newtcolorbox{routemap}{
  enhanced, breakable,
  colback=edu-green-bg,
  colframe=edu-green!30,
  boxrule=0pt,
  borderline west={2.5pt}{0pt}{edu-green},
  left=3mm, right=3mm, top=3mm, bottom=3mm,
}

% 关键想法 — 蓝色左边线
\newtcolorbox{keyideabox}{
  enhanced, breakable,
  colback=edu-blue-bg,
  colframe=edu-blue!30,
  boxrule=0pt,
  borderline west={2.5pt}{0pt}{edu-blue},
  left=3mm, right=3mm, top=3mm, bottom=3mm,
}

% 易错提醒 — 红色左边线
\newtcolorbox{mistakebox}{
  enhanced, breakable,
  colback=edu-red-bg,
  colframe=edu-red!20,
  boxrule=0pt,
  borderline west={3pt}{0pt}{edu-red},
  left=3mm, right=3mm, top=2.5mm, bottom=2.5mm,
}

% 提示 — 灰色虚线左边线
\newtcolorbox{hintbox}{
  enhanced, breakable,
  colback=edu-gray-bg,
  colframe=edu-gray!20,
  boxrule=0pt,
  borderline west={2pt}{0pt}{edu-gray, dashed},
  left=3mm, right=3mm, top=2mm, bottom=2mm,
}

% 思考题 — 紫色虚线左边线
\newtcolorbox{thinkbox}{
  enhanced, breakable,
  colback=edu-purple-bg,
  colframe=edu-purple!20,
  boxrule=0pt,
  borderline west={2pt}{0pt}{edu-purple, dashed},
  left=2.5mm, right=2.5mm, top=2.5mm, bottom=2.5mm,
}

% 教师备注 — 灰色左边线，小字
\newtcolorbox{teachernote}{
  enhanced, breakable,
  colback=edu-gray-bg,
  colframe=edu-gray!15,
  boxrule=0pt,
  borderline west={2pt}{0pt}{edu-gray!60},
  left=2.5mm, right=2.5mm, top=2.5mm, bottom=2.5mm,
  fontupper=\small,
  colupper=black!60,
}

% 训练目标 — 橙色左边线，小字
\newtcolorbox{traininggoal}{
  enhanced, breakable,
  colback=edu-orange-bg,
  colframe=edu-orange!15,
  boxrule=0pt,
  borderline west={1.5pt}{0pt}{edu-orange},
  left=2mm, right=2mm, top=1mm, bottom=1mm,
  fontupper=\small,
}

% 答题区 — 无边框，纯留白
\newcommand{\answerarea}[1][40mm]{%
  \vspace{2mm}%
  \begin{tcolorbox}[
    enhanced, breakable,
    colback=white,
    colframe=white,
    boxrule=0pt,
    height=#1,
    valign=top,
    bottom=0mm,
    after skip=0mm,
  ]
  \end{tcolorbox}%
}

% ---------- 字体设置 ----------
% exam-zh (ctexbook) already sets CJK fonts via ctex.
% Only override if specific fonts are needed.
% macOS: Songti SC, STHeiti, STFangsong
% Linux: SimSun, SimHei, FangSong

% ---------- 页面设置 ----------
% exam-zh (ctexbook) already loads geometry.
% Override margins if needed:
% \geometry{top=16mm, bottom=16mm, left=14mm, right=14mm}

\examsetup{
  question/show-answer = false,
  fillin/show-answer = false,
  solution/show-solution = hide,
}

% ---------- 教师版解答样式 ----------
\newtcolorbox{solutionblock}{
  enhanced, breakable,
  colback=edu-blue-bg,
  colframe=edu-blue!25,
  boxrule=0pt,
  borderline west={2pt}{0pt}{edu-blue!50},
  arc=0pt,
  left=3mm, right=3mm, top=2mm, bottom=2mm,
  before skip=2mm,
  after skip=3mm,
  fontupper=\small,
}

% 解答步骤编号
\newcounter{solstep}
\newcommand{\solstep}[2]{%
  \stepcounter{solstep}%
  \par\medskip
  \noindent\hangindent=6mm
  {\color{edu-blue}\bfseries\textcircled{\scriptsize\thesolstep}\enspace #1}\enspace #2\par
}

\begin{document}

\section*{通过交点、面积、距离等条件求点坐标}

\section*{一、选择题（每题 6 分，共 36 分）}



\begin{question}[points=6]
  直线 $y=2x+b$ 过点 $(3,8)$，则 $b$ 的值为
  \begin{choices}
    \item $2$
    \item $-2$
    \item $14$
    \item $22$
  \end{choices}
\end{question}



\begin{question}[points=6]
  直线 $y=2x+1$ 与 $y=-x+4$ 的交点坐标为
  \begin{choices}
    \item $(1, 3)$
    \item $(3, 1)$
    \item $(2, 5)$
    \item $(5, 2)$
  \end{choices}
\end{question}



\begin{question}[points=6]
  直线 $y=-3x+6$ 与 $x$ 轴交于 $A$，过原点作该直线的平行线与 $y$ 轴交于 $B$，则
  \begin{choices}
    \item $A(2,0)$，$B$ 与原点重合
    \item $A(2,0)$，$B(0,6)$
    \item $A(0,6)$，$B(2,0)$
    \item $A(2,0)$，平行线过原点即 $y=-3x$，与 $y$ 轴交于 $(0,0)$
  \end{choices}
\end{question}



\begin{question}[points=6]
  $O(0,0)$、$A(4,0)$，$P$ 在直线 $y=x$ 上且 $\triangle OAP$ 面积为 $6$，则 $P$ 的纵坐标为
  \begin{choices}
    \item $2$
    \item $4$ 或 $-4$
    \item $3$ 或 $-3$
    \item $6$
  \end{choices}
\end{question}



\begin{question}[points=6]
  直线 $y=-2x+6$ 与坐标轴围成的 $\triangle AOB$ 中，点 $C$ 在 $OA$ 上使 $\triangle OBC$ 面积为 $\triangle AOB$ 的 $\dfrac{1}{3}$，则 $OC$ 的长为
  \begin{choices}
    \item $1$
    \item $2$
    \item $3$
    \item $6$
  \end{choices}
\end{question}



\begin{question}[points=6]
  $A(0,3)$，$B$ 在第一象限且 $AB \perp y$ 轴，$AB=5$，则 $B$ 的坐标为
  \begin{choices}
    \item $(5, 3)$
    \item $(3, 5)$
    \item $(5, 8)$
    \item $(3, 0)$
  \end{choices}
\end{question}


\section*{二、填空题（每题 6 分，共 24 分）}



\begin{question}[points=6]
  直线 $y=-x+4$ 与 $x$ 轴、$y$ 轴分别交于 $A$、$B$，则 $\triangle AOB$ 的面积为 \rule{3cm}{0.4pt}。
  \fillin
\end{question}



\begin{question}[points=6]
  点 $P$ 在直线 $y=2x-4$ 上，$\triangle OPA$ 中 $O(0,0)$、$A(3,0)$，若面积为 $5$，则 $P$ 的坐标为 \rule{3cm}{0.4pt}。
  \fillin
\end{question}



\begin{question}[points=6]
  点 $A(1,2)$、$B(4,6)$，则线段 $AB$ 的长为 \rule{3cm}{0.4pt}。
  \fillin
\end{question}



\begin{question}[points=6]
  点 $A(-3,5)$ 关于 $y$ 轴的对称点为 $A'$，则 $A'$ 的坐标为 \rule{3cm}{0.4pt}。
  \fillin
\end{question}


\section*{三、解答题（共 40 分）}


\needspace{8\baselineskip}

\begin{problem}[points=12]
  直线 $y = -\dfrac{4}{3}x + 4$ 与 $x$ 轴、$y$ 轴分别交于 $A$、$B$。

(1) 求 $A$、$B$ 的坐标；
(2) 点 $P$ 在直线 $AB$ 上，$\triangle OAP$ 的面积为 $4$，求 $P$ 的坐标。
  \answerarea[60mm]
\end{problem}


\needspace{8\baselineskip}

\begin{problem}[points=13]
  直线 $l_1$ 过 $A(0,3)$、$B(2,0)$，直线 $l_2$ 过 $B(2,0)$、$C(4,4)$。

(1) 求 $l_1$ 和 $l_2$ 的解析式；
(2) 过 $A$ 作 $l_2$ 的平行线交 $x$ 轴于 $D$，求 $D$ 的坐标。
  \answerarea[60mm]
\end{problem}


\needspace{8\baselineskip}

\begin{problem}[points=15]
  $A(0,4)$，$B(6,0)$，$C$ 在 $x$ 轴上（$C$ 在 $O$、$B$ 之间），$\triangle ABC$ 中 $AC = BC$。

(1) 求 $C$ 的坐标；
(2) 求 $AC$ 所在直线的解析式。
  \answerarea[65mm]
\end{problem}



\end{document}
